% Generated draft: requires mathematical and source review.
\documentclass[11pt]{article}
\usepackage[margin=25mm]{geometry}
\usepackage{fontspec}
\setmainfont{Latin Modern Roman}
\usepackage{amsmath,amssymb,mathtools,unicode-math}
\setmathfont{Latin Modern Math}
\usepackage{enumitem,needspace,xurl,hyperref}
\setlength{\parindent}{0pt}
\setlength{\parskip}{5pt}
\setlength{\emergencystretch}{3em}
\newcommand{\N}{\mathbb{N}}
\newcommand{\Z}{\mathbb{Z}}
\newcommand{\Q}{\mathbb{Q}}
\newcommand{\R}{\mathbb{R}}
\newcommand{\C}{\mathbb{C}}
\newcommand{\F}{\mathbb{F}}
\newcommand{\A}{\mathbb{A}}
\newcommand{\PP}{\mathbb{P}}
\newcommand{\E}{\mathbb{E}}
\newcommand{\eps}{\varepsilon}
\newcommand{\dd}{\,\mathrm{d}}
\begin{document}
\section*{UnsolvedMath ID: 96}
\textbf{Source problem code:} catalogue entry
\subsection{Definitions and mathematical statement}
The intended Ulam sequence $U=\{u_1,u_2,\ldots\}\subset\N$ starts with $u_1=1$, $u_2=2$ and has the displayed initial terms
\[
1,2,3,4,6,8,11,13,16,18,26,28,36,\ldots.
\]
For $n\geq2$ and a positive integer $m$, define
\[
r_n(m)=\bigl|\{(i,j):1\leq i<j\leq n,\ u_i+u_j=m\}\bigr|.
\]
The usual increasing-sequence interpretation of the source recursion is
\[
u_{n+1}=\min\{m\in\N:m>u_n,\ r_n(m)=1\}.
\]
Representations use two distinct earlier terms and are counted without ordering by requiring $i<j$. The condition $m>u_n$ is essential to this intended interpretation: the literal supplied sentence omits it, although its displayed strictly increasing sequence indicates it. Without that condition the recursion can select a previously used term. This qualification must be checked against the intended source.

Write $A(N)=|U\cap\{1,\ldots,N\}|$ for positive integers $N$. One precise reading of ``positive density'' asks whether
\[
\lim_{N\to\infty}\frac{A(N)}{N}=d\quad\hbox{exists with }d>0.
\]
A different, weaker reading asks only whether
\[
\liminf_{N\to\infty}\frac{A(N)}{N}>0.
\]
The supplied question does not distinguish natural density from positive lower density, so these are recorded as distinct possible targets rather than identified.

The additional question asks for an explanation of the sequence's curious Fourier properties. For example, finite Fourier sums can be written
\[
F_N(\theta)=\sum_{\substack{u\in U\\u\leq N}}\exp(2\pi i\theta u),\qquad \theta\in\R/\Z,
\]
but the source specifies no particular frequencies, normalization, limiting statement, or estimate to prove. This is therefore an exploratory part of the problem, not a quantified Fourier conjecture. The supplied background reports numerical density near $0.07$ and apparently quasi-periodic Fourier structure; neither observation is stated as a theorem. Its triage dated 2026-08-17 labels positive density and a structural explanation open, excludes speculative/AI notes as results, and reports that no rigorous positive-density theorem was found in that review.
\subsection{Short English statement}
For the standard increasing Ulam sequence, which repeatedly adds the next integer having exactly one representation as a sum of two distinct previous terms, prove positive density and explain the observed Fourier structure. The source leaves the exact density notion and Fourier target unspecified, and its written recursion omits the required increasing condition.
\subsection{Sources}
\begin{itemize}
\item Ben Green problem collection, catalogue entry, checked 2026-08-17 according to the supplied literature triage: \url{https://www.unsolvedmath.com/}. The record supplies this as its reference for the open framing and numerical context.
\item Local registry data/unsolvedmath/problems.json, permanent catalogue entry, source code catalogue entry, ``Ulam's Sequence,'' attributed to Stanis\l aw Ulam, in Ben Green's 100 Open Problems; complete statement, background, and literature triage dated 2026-08-17. Dataset attribution: Ulam AI and the UnsolvedMath Contributors.
\end{itemize}
\end{document}
